\subsection{Roadmap M26--M30}

\begin{itemize}[leftmargin=1.5em]
\item \textbf{M26.} Fused exact hot/cold runtime. Стадия 1 --- allocation-free
Torch path (\texttt{full\_weight\_hot\_v2}); landed. На same-encoding
\texttt{l54\_q\_gu} 16w даёт \texttt{PPL = 2.7996} (без изменений), пиковая
VRAM \texttt{46385.9} MB (без изменений), throughput \texttt{971.28} tok/s
против \texttt{962.06} у \texttt{full\_weight\_hot} и \texttt{943.29} у
\texttt{full\_weight\_fast} ($+0.96\%$ и $+2.97\%$ соответственно). Целевой
рубеж $\ge 980$ tok/s на этом узком слайсе пока не достигнут. На
\texttt{l0\_qkv\_gu} 16w (патчатся \texttt{q,k,v,gate,up} слоя 0): тот же
PPL \texttt{32.2524}, throughput \texttt{961.86} tok/s против \texttt{948.17}
и \texttt{947.92}; одновременно зафиксирована регрессия по VRAM в hot-кэше
(\texttt{88.9} GB против \texttt{47.2} GB у fast), которую стадия 2 должна
переписать. Была проверена и более дешёвая промежуточная ветка
\texttt{triton\_hot\_cold\_persistent}: exact tile-local persistent palette +
hotprefix Triton kernel. Она оказалась отрицательным результатом: на honest
same-encoding \texttt{l54\_q\_gu} 16w первый run дал \texttt{PPL = NaN} и лишь
\texttt{188.77} tok/s при peak \texttt{47567.1} MB; последующая правка,
строившая persistent plan из того же pre-row-scale tensor, что и
\texttt{full\_weight\_fast}, улучшила synthetic parity, но полный harness всё
равно упёрся в CUDA gather out-of-bounds. Следовательно, стадия 2 должна идти
не через repack всей матрицы в local palettes, а через настоящий stage-local
RVQ fused hot/cold matmul. Первый native probe в этом направлении тоже уже
сделан: \texttt{stage\_local\_hot\_cold} работает поверх родных
\texttt{stage\_ids + full\_codebooks + hot\_lut + hot\_codebooks}. На one-hot
input для \texttt{l54.gate\_proj} / \texttt{l54.up\_proj} он проходит narrow
correctness gate (max abs diff порядка $10^{-4}$ против
\texttt{full\_weight\_hot\_v2}), то есть сама hot/cold логика на native RVQ
surface собрана правильно. Но это всё ещё scaffold, а не победа: на
реалистичном $m = 2048$ microbench он медленнее \texttt{full\_weight\_hot\_v2}
(примерно \texttt{58.8 -> 90.0} ms на \texttt{gate\_proj}), а на honest
same-encoding \texttt{l54\_q\_gu} 16w даёт \texttt{PPL = NaN}, лишь
\texttt{258.23} tok/s и peak \texttt{46575.8} MB. Попытка B1 убрать LUT на
уровне API тоже отрицательна: \texttt{stage\_local\_hot\_cold\_b1} сохраняет
one-hot parity, но realistic gate проваливает ещё жёстче
(\texttt{58.81 -> 569.43} ms на \texttt{gate\_proj}). B2 исправил local
compute shape: staged hot palette теперь даёт one-hot max abs diff
\texttt{0.001953125}, \texttt{has\_nan = False} и примерно
\texttt{78.84 -> 68.36} ms ($1.153\times$) против
\texttt{full\_weight\_hot\_v2}. Но probe всё равно отрицательный: per-stage
hot-hit rate лишь около \texttt{0.075}, то есть на порядок ниже intended
порога \texttt{0.72}. B3 проверил более сильную гипотезу напрямую: оставить B2
kernel, но заменить policy выбора hot IDs. Для этого runtime был доведён до
реально модульного состояния, включая поддержку non-power-of-2
\texttt{hot\_topk} через padding staged hot palette. Однако selection-only fix
тоже не спасает stage 2: на authoritative \texttt{l54.gate/up}, $m = 2048$,
режимы \texttt{topk = 48} и \texttt{64} сохраняют ту же локальную численность
(one-hot max abs diff \texttt{0.001953125}, \texttt{has\_nan = False}), но дают
лишь \texttt{0.170--0.222} mean hit rate и около \texttt{0.224x} от
\texttt{full\_weight\_hot\_v2}. Более того, count-based upper bound показывает,
что даже чистый stage-local \texttt{count} selection даёт лишь
\texttt{~0.312} hit mean при \texttt{topk = 64}; при \texttt{topk = 128}
coverage поднимается только до \texttt{~0.578}, но speed падает примерно до
\texttt{0.133x}. Значит, текущий M26-вопрос формулируется уже жёстче: не просто
"какая hot-selection policy лучше", а "поддерживает ли текущее Block-RVQ
encoding вообще маленький stage-local hot palette на этих MLP-слоях". Целевой
рубеж стадии 2 остаётся $\ge 1100$ tok/s при том же качестве, но ни B0--B3 его
пока не приблизили честным образом.

Следующая компиляторная гипотеза была ещё более прямой: отказаться от словаря
"hot/cold" и проверить \emph{Route Register File + spilling}, то есть явный
register-allocation pass над токенами $(stage,id)$. Но честный Step~1a для
этой идеи был поставлен как offline gate до любого kernel work: сначала
собрать activation-weighted influence, occurrence counts и structural
interference per stage, а уже потом смотреть, есть ли у register allocator
вообще какой-либо рычаг. И этот gate дал ещё более жёсткий отрицательный
результат. На \texttt{l54.gate\_proj} / \texttt{l54.up\_proj} при tile-size,
равном \texttt{block\_n = 256} B2-kernel'а, каждый stage имеет
\texttt{avg\_id\_tile\_occupancy = 1.000}: все codebook id оказываются живыми
в каждом tile. Следовательно, interference graph практически полносвязный, а
linear-scan allocator с приоритетом
\texttt{Influence $\times$ (1 - avg\_interference)} вырождается в
\texttt{topk-by-influence}. И это не даёт даже локального выигрыша: mean
hit-rate на \texttt{12} stages составляет лишь \texttt{0.293} при
\texttt{cap = 64} и \texttt{0.560} при \texttt{cap = 128}, что хуже или равно
простому \texttt{topk-count} baseline (\texttt{0.311} и \texttt{0.577}).
Поэтому build полноценного RRF runtime branch на текущем encoding был
остановлен ещё до Step~1b. Этот отрицательный результат важен концептуально:
он показывает, что даже честная compiler-style register allocation не помогает,
если stage vocabulary уже structurally full на уровне tile. Значит, следующий
M26/M27 шаг должен менять либо само encoding-пространство, либо compute
scheduling, а не добавлять ещё один selection layer поверх тех же токенов.

Следующая гипотеза проверяла уже именно tile-local compilation: вместо одного
register file на всю stage собрать для каждого тайла веса размера
\texttt{64 x 256} собственную маленькую palette и выбирать в неё
\texttt{topk = 32/48/64} id по tile-local activation-weighted influence.
Это и был честный offline gate для \emph{Per-Tile Dynamic Palette (PTDP)}.
Если бы проблема RRF была только в stage-global scope, такой шаг должен был бы
резко поднять tile-local coverage. Но этого не произошло. На
\texttt{l54.gate\_proj} / \texttt{l54.up\_proj} mean tile count-hit across
stages составляет лишь \texttt{0.286/0.393/0.488} и
\texttt{0.289/0.394/0.490} для \texttt{topk = 32/48/64}; лучший stage mean во
всём прогоне --- только \texttt{0.494}, а доля тайлов с
\texttt{count\_hit >= 0.75} равна \texttt{0.000} для каждого проверенного
\texttt{topk}. Поэтому PTDP Step~2b, то есть runtime integration поверх B2
kernel'а, честно даже не запускался. Этот результат уточняет предыдущий вывод:
дело уже не только в полносвязном interference graph на stage-level, но и в
том, что текущий Block-RVQ encoding не создаёт достаточно концентрированной
лексики даже внутри tile. Значит, сам вопрос для M26/M27 теперь звучит ещё
жёстче: можно ли вообще добиться маленького cached vocabulary на этих MLP-слоях
без смены representation, а не только без смены selection policy.

Следующий шаг проверял уже сам язык, а не selector: \emph{Finer-Grained Route
Language (FGRL)}. Идея состояла в том, чтобы поднять суммарную residual depth
с текущих \texttt{12} stages до \texttt{20}, но одновременно уменьшить
stage-local alphabet с \texttt{256} до \texttt{80}. В терминах этого repo это
означало не буквальный \texttt{num\_stages = 20} на split, а честную
реализацию total-20 через \texttt{stages\_per\_split = (5,5,5,5)} при
\texttt{product\_splits = 4}. На двухслойном gate для
\texttt{l54.gate\_proj} / \texttt{l54.up\_proj} FGRL действительно улучшил
качество: average layer relMSE упал с \texttt{0.0390} до \texttt{0.0172}, а
stable \texttt{4}-window \texttt{full\_weight\_fast} eval улучшил PPL с
\texttt{2.4054} до \texttt{2.3928}. Tile \texttt{top64} share тоже вырос
сильно: с \texttt{0.495} до \texttt{0.909}. Но решающий механизм снова не
сработал. Нормализованная tile occupancy не падает, а растёт:
\texttt{0.853 -> 0.996}. То есть почти весь новый \texttt{80}-id vocabulary
живёт в каждом tile, а высокий \texttt{top64} share получается в основном
потому, что \texttt{64} уже почти покрывает весь alphabet, а не потому, что
язык стал genuinely sparse. Одновременно растёт storage cost
(\texttt{3.004 -> 5.003} bits/weight), throughput падает
(\texttt{1483.29 -> 1309.01} tok/s), а eval peak VRAM чуть увеличивается
(\texttt{46486.8 -> 46598.5} MB). Поэтому FGRL Step~3b как runtime branch тоже
не был открыт. Этот результат особенно важен: он показывает, что even true
language redesign может улучшать quality, но всё ещё не решать аппаратную
задачу small cached vocabulary, если новая лексика остаётся почти fully live
на уровне tile.

Последний на данный момент шаг проверял уже не новую palette policy и не новый
alphabet, а сам syntax layer языка: \emph{WAL with Structured Syntax (WAL-SS)}.
В первой, максимально консервативной формулировке это был строго lossless
macro layer поверх текущих \texttt{12 x 256} stage-id программ. Для
\texttt{l54.gate\_proj} / \texttt{l54.up\_proj} были выделены до
\texttt{128} frequent subsequences длиной \texttt{3-5}, после чего каждый
block-program переписывался как \texttt{macro calls + literal args} с точным
decode обратно в исходный stream. Correctness gate здесь проходит идеально:
stage IDs совпадают exactly, \texttt{recon\_rel\_mse = 0.0}, а stable
\texttt{4}-window \texttt{full\_weight\_fast} eval идентичен current packed
path (\texttt{PPL = 2.4081}, \texttt{1477.69 tok/s}, peak
\texttt{46486.8} MB). Но структурный выигрыш почти отсутствует. На каждом из
двух слоёв имеется около \texttt{7,340,032} block-programs, и почти каждый из
них уникален (\texttt{7,340,018} / \texttt{7,340,015} unique full programs).
Поэтому syntax layer почти не сжимает программу: average program length падает
лишь с \texttt{12} до \texttt{11.9985} / \texttt{11.9987}, compression ratio
остаётся \texttt{0.999875} / \texttt{0.999889}, а macro token coverage равен
всего \texttt{0.000176} / \texttt{0.000158}. Это важный отрицательный
результат уже для WAL как языка: даже идеальный lossless macro wrapper не
создаёт usable structured syntax, если сам stage-program stream почти полностью
уникален. Значит, следующий WAL-level шаг, если он вообще будет syntax-based,
должен строить не просто dictionary compression поверх current encoding, а
настоящую повторяемую higher-order structure.

Именно эту более сильную syntax-гипотезу затем проверил первый шаг
\emph{WAL with Hierarchical Programs (WAL-HP)}. В первой честной формулировке
он тоже оставался строго exact: были взяты уже четыре MLP-слоя
(\texttt{l53/l54 gate/up}), после чего глобально, across layers, были
выделены \texttt{256} shared subsequences длиной \texttt{3-6} и каждая
block-program переписана как \texttt{CALL shared\_subroutine} или
\texttt{LITERAL}. Correctness gate снова проходит идеально:
\texttt{exact\_stage\_ids = True}, \texttt{recon\_rel\_mse = 0.0}, а stable
\texttt{4}-window \texttt{full\_weight\_fast} gate идентичен current packed
path на этом четырёхслойном slice (\texttt{PPL = 2.4258},
\texttt{1216.91 tok/s}, peak \texttt{47530.4} MB). Более того, некоторый
настоящий cross-layer reuse уже появляется: \texttt{136 / 256} selected
subroutines используются хотя бы в двух target'ах, а strongest motif
(\texttt{[128,128,128]}) встречается во всех четырёх. Но и этого снова
недостаточно для нового ISA baseline. Global \texttt{call\_coverage} равен
лишь \texttt{0.000166}, а average program length остаётся практически flat:
\texttt{11.9973 / 11.9989 / 11.9989 / 11.9992} при raw length \texttt{12}.
То есть even shared hierarchical subroutines пока существуют лишь как trace
motifs, а не как meaningful high-level program surface. Поэтому WAL-HP
Step~5b честно не запускался. Если двигаться дальше по syntax-линии, то уже
не к ещё одному exact wrapper, а к более сильной форме структуры:
parameterized или approximate subroutines, learned templates, либо encoder,
который сам создаёт повторяемые higher-order blocks.

Следующий probe проверил уже именно learned-template вариант этой идеи:
\emph{WAL with Learned Route Templates (WAL-LRT)}. На том же двухслойном
slice \texttt{l54.gate/up} был обучен банк из \texttt{256} categorical
full-program templates длины \texttt{12}, после чего отдельно были измерены
две поверхности: exact \texttt{TEMPLATE + literal corrections} и
\texttt{template\_only} approximate encoding. Это первый syntax-level шаг,
который действительно создаёт немного больше искусственной структуры, чем
WAL-SS/WAL-HP: average program length падает до
\texttt{11.7495 / 11.7494} при raw length \texttt{12}, а template token
coverage поднимается до \texttt{0.04137 / 0.04130}. Exact wrapper, как и
должно быть, сохраняет качество полностью: stable \texttt{4}-window
\texttt{full\_weight\_fast} gate остаётся идентичен current packed path
(\texttt{PPL = 2.4081}, \texttt{1479.74 tok/s}, peak \texttt{46486.8} MB).
Но learned templates сами по себе ещё слишком слабы как semantics-bearing
surface. Nearest template совпадает лишь примерно по \texttt{10.4\%} токенов
(average Hamming distance около \texttt{10.75 / 12}), reconstruction relMSE
взлетает до примерно \texttt{1.84 / 1.84}, а \texttt{template\_only} PPL
ухудшается до \texttt{2.9375}. Поэтому WAL-LRT Step~6b честно не запускался:
первый learned discrete template bank уже создаёт немного синтаксиса, но ещё
не создаёт usable high-level ISA без большой correction tail. Следующий
template-level шаг, если он будет, должен делать templates более выразительными
самими по себе: parameterized templates, weight-space templates или continuous
corrections.

Следующий шаг сделал syntax-гипотезу ещё буквальнее:
\emph{WAL as a Formal Grammar (WAL-FG)}. На том же двухслойном slice
\texttt{l54.gate/up} была задана shallow grammar из \texttt{15} правил:
\texttt{ROOT -> LEFT RIGHT}, затем \texttt{LEFT/RIGHT -> SLOT SLOT}, а у
каждого из четырёх phrase-slots длины \texttt{3} есть по \texttt{3} learned
productions. Это даёт каждой block-program длины \texttt{12} формальное дерево
глубины \texttt{4} с \texttt{7} nonterminal nodes. Были отдельно измерены две
поверхности: exact parse с \texttt{RULE calls + literal corrections} и
\texttt{grammar\_only} approximate encoding без corrections. Exact parse
действительно создаёт formal grammar surface, но почти не меняет экономику
языка: average program length остаётся \texttt{11.9984 / 11.9988} при raw
length \texttt{12}, а rule token coverage равен лишь
\texttt{0.000272 / 0.000203}. Approximate grammar оказывается чуть менее
катастрофической, чем template-only в WAL-LRT, но всё ещё недопустимой по
quality bar: reconstruction relMSE поднимается до
\texttt{1.2928 / 1.3448} (eval-layer relMSE около \texttt{1.31}), а stable
\texttt{4}-window \texttt{full\_weight\_fast} gate ухудшается с
\texttt{PPL = 2.4081} до \texttt{2.6866} при практически той же скорости и
той же peak VRAM. Поэтому WAL-FG Step~7a тоже является честным отрицательным
результатом: наличие formal parse tree ещё не означает наличие usable weight
grammar. Следующий grammar-level шаг, если он вообще будет, должен искать не
просто дерево, а productions, которые действительно несут semantics.

Следующая гипотеза подняла ещё один уровень формальности:
\emph{WAL with Route Types and Type System (WAL-TS)}. На том же slice
\texttt{l54.gate/up} был задан простой typed layer из шести atomic types:
\texttt{MLP\_GATE/UP x COMMON/MIXED/OUTLIER}. Каждый block-program получает
тип по простой эвристике среднего token surprisal, после чего для каждого
локального типа обучается собственный phrase bank и grammar уже строится как
\texttt{TYPE + RULE calls + literal corrections}. Этот шаг действительно
создаёт явное label-space: в каждом из двух слоёв активны все три локальных
типа, common/outlier блоки слегка отделяются друг от друга, а их доли устойчиво
лежат около \texttt{48.9\% / 40.8\% / 10.3\%}. Но typed grammar всё равно
остаётся почти полностью literal: average program length равен лишь
\texttt{11.99815 / 11.99848} при raw length \texttt{12}, а typed rule coverage
составляет только \texttt{0.000308 / 0.000253}. Более того, approximate
\texttt{typed\_only} surface оказывается не лучше, а хуже WAL-FG на этом
slice: reconstruction relMSE поднимается до \texttt{1.6587 / 1.9014}
(eval-layer relMSE около \texttt{1.75}), а stable \texttt{4}-window gate
ухудшается с \texttt{PPL = 2.4081} до \texttt{2.7994} при практически той же
скорости и той же peak VRAM. Поэтому WAL-TS Step~8a тоже является честным
отрицательным результатом: наличие typed labels ещё не означает наличие
semantics-bearing typed assembly. Если продолжать эту линию, нужны типы,
которые действительно ограничивают и объясняют поведение программы, а не
только переименовывают rarity bands.

Следующая гипотеза подняла syntax-layer до максимально буквальной формы:

\emph{WAL as Classic Assembly (WAL-ASM)}. На том же slice
\texttt{l54.gate/up} был собран explicit assembly surface с
\texttt{LABEL/CALL/JMP/MACRO/RET}. Для этого были обучены \texttt{64}
full-program subroutines длины \texttt{12}, каждая из которых становится
именованной подпрограммой \texttt{gate\_sub\_xxx / up\_sub\_xxx}; затем их
тела дополнительно macroize-ятся \texttt{32} macros длины \texttt{3--5}, а
между соседними \texttt{CALL}-сайтами строится dominant jump-table. Были
отдельно измерены exact surface с \texttt{CALL + literal corrections} и
approximate \texttt{asm\_only} surface без corrections. Exact WAL-ASM, как и
должно быть, полностью сохраняет quality gate current packed path
(\texttt{PPL = 2.4081}, \texttt{1481.11 tok/s}, peak \texttt{46486.8} MB),
но control-flow остаётся только trace-level: average program length равен
лишь \texttt{11.9288 / 11.9309} при raw length \texttt{12}, average calls ---
\texttt{0.0702 / 0.0681}, average jumps --- всего
\texttt{0.000177 / 0.000174}, call-token coverage --- только
\texttt{0.0118 / 0.0114}, а macro-body coverage ---
\texttt{0.0717 / 0.0751}. Approximate \texttt{asm\_only} surface ещё хуже с
точки зрения semantics: token match составляет лишь
\texttt{0.08595 / 0.08580}, reconstruction relMSE поднимается до
\texttt{1.8789 / 1.8749} (eval-layer relMSE около \texttt{1.85}), а stable
\texttt{4}-window gate ухудшается до \texttt{PPL = 2.9601} при практически
той же скорости и той же peak VRAM. Поэтому WAL-ASM Step~9a тоже является
честным отрицательным результатом: наличие классических assembly-механизмов
ещё не означает наличие semantics-bearing weight program. Если продолжать
эту линию, нужны не просто labels и jumps сами по себе, а representation или
assembly semantics, где \texttt{CALL/JMP} покрывают заметную часть поведения,
а не остаются следовым wrapper поверх nearest templates.

Следующая гипотеза впервые отказалась чинить syntax-layer сверху и попробовала
спроектировать саму ISA:

\emph{WAL with Learned Discrete ISA + Semantic Constraints (WAL-LDI)}.
На том же slice \texttt{l54.gate/up} был построен двухуровневый язык:
сначала joint feature+sequence clustering выделяет \texttt{4} learned semantic
families на слой, затем внутри каждой family обучается по \texttt{4} low-level
atoms для каждого из \texttt{4} phrase-slots длины \texttt{3}. Итого получается
\texttt{68} инструкций на слой: \texttt{4} high-level family opcodes и
\texttt{64} family-conditioned low-level instructions. Это уже не просто
surface wrapper. Впервые появляется реальная semantic partition: все
\texttt{4} family активны, family entropy держится около \texttt{1.14}, а
маленький \texttt{COMMON\_CORE} cluster отделяется очень сильно
(margin около \texttt{6.36} на \texttt{gate} и \texttt{6.11} на \texttt{up})
и даёт лучший local approximate match (около \texttt{0.113 / 0.109}).
Но usable ISA из этого пока не получается. Exact program surface даже длиннее
raw: average program length равен \texttt{12.9969 / 12.9971} при raw length
\texttt{12}, потому что каждый block всегда несёт high-level family header,
а low-level atoms почти не срабатывают
(\texttt{avg low-level calls = 0.00312 / 0.00289}, token coverage только
\texttt{0.000521 / 0.000483}). Approximate \texttt{ldi\_only} surface уже
лучше, чем \texttt{asm\_only}, но всё ещё далека от sacred quality bar:
reconstruction relMSE поднимается до \texttt{1.7697 / 1.8243}
(eval-layer relMSE около \texttt{1.77}), а stable \texttt{4}-window gate
ухудшается до \texttt{PPL = 2.8318} при практически той же скорости и той же
peak VRAM. Поэтому WAL-LDI Step~10a --- это не победа, а первый частично
положительный сигнал в правильном направлении: semantic families действительно
можно выделить, но если делать это post-hoc поверх уже почти случайного
stage-stream, язык всё ещё не становится compact semantics-bearing ISA.
Следующий viable LDI-level шаг должен встраивать semantic constraints в сам
encoder/ISA learning, а не только кластеризовать готовую программу после факта.

Следующий шаг проверял именно эту более сильную гипотезу:

\emph{WAL with End-to-End Learned ISA + Semantic Encoder (WAL-E2E)}.
На том же slice \texttt{l54.gate/up} был добавлен маленький semantic encoder,
который jointly предсказывает high-level family ID и family-conditioned
low-level atom IDs по \texttt{4} phrase-slots длины \texttt{3}, а сами atoms
обучаются одновременно с encoder'ом. Objective объединяет token reconstruction,
semantic neighborhood consistency, family cohesion и anti-collapse
regularization. Это действительно исправляет один важный structural defect:
в authoritative run все \texttt{4} family остаются активны на обоих слоях,
family entropy составляет \texttt{1.147} на \texttt{gate} и \texttt{0.802}
на \texttt{up}, а coherent \texttt{COMMON\_CORE} island не исчезает
(около \texttt{1.97\%} блоков и margin \texttt{6.34} на \texttt{gate};
около \texttt{7.43\%} и margin \texttt{2.76} на \texttt{up}).
Но usable ISA из этого всё равно не получается. Exact program surface остаётся
ещё длиннее, чем в WAL-LDI: average program length равен
\texttt{12.9988 / 12.9988} при raw length \texttt{12}; low-level atoms почти
не используются (\texttt{avg low-level calls = 0.00119 / 0.00118}, token
coverage лишь \texttt{0.000199 / 0.000197}). Approximate \texttt{e2e\_only}
surface по quality тоже не проходит sacred bar:
eval-layer relMSE составляет около \texttt{1.70}, а stable \texttt{4}-window
gate ухудшается до \texttt{PPL = 2.8847}, что даже хуже первого WAL-LDI
baseline. Поэтому WAL-E2E Step~11a --- тоже честный отрицательный результат,
но уже более содержательный: недостаточно просто перенести semantic ISA внутрь
encoder loop, если objective всё ещё учит в основном token reconstruction на
готовом stage language. Следующий viable E2E-level шаг должен оптимизировать
не только token CE, а напрямую weight/block reconstruction или layer-output
semantics вместе с explicit program-cost / correction-budget loss.

\emph{WAL with Direct Block Reconstruction + Explicit Program Cost (WAL-DR)}.
Следующий шаг уже меняет сам objective. На том же slice \texttt{l54.gate/up}
сохраняется semantic encoder и family-conditioned atom bank, но основной loss
теперь оптимизирует direct final-block reconstruction под текущим Block-RVQ
surface, а отдельный program-cost proxy штрафует correction-heavy exact
program. Это даёт первый содержательный сигнал, что bottleneck действительно
сидел в objective: все \texttt{4} family остаются активны (family entropy
\texttt{1.1188} на \texttt{gate} и \texttt{1.2486} на \texttt{up}), а
approximate \texttt{dr\_only} surface становится заметно лучше, чем и у
WAL-LDI, и у WAL-E2E: eval-layer relMSE падает до примерно \texttt{1.54}, а
stable \texttt{4}-window gate ухудшается лишь до \texttt{PPL = 2.7325} вместо
\texttt{2.8318} у \texttt{ldi\_only} и \texttt{2.8847} у
\texttt{e2e\_only}. Но compressible ISA из этого всё ещё не получается.
Exact program остаётся длиннее raw: average program length равен
\texttt{12.9984 / 12.9988} при raw length \texttt{12}, low-level atoms всё
ещё почти не используются (\texttt{avg low-level calls = 0.00156 / 0.00117},
token coverage \texttt{0.000260 / 0.000195}). Поэтому WAL-DR Step~12a --- это
не новая baseline ISA, а первый положительный objective-level correction:
quality-side заметно улучшается, но exact program всё ещё остаётся почти
полностью correction-dominated. Следующий viable шаг должен привязывать cost
ещё ближе к downstream behavior --- например, к layer-output reconstruction
и/или к более явному выбору между high-level calls и literal corrections.

\emph{WAL with Layer-Output Reconstruction + Strong Program Cost (WAL-LO)}.
Следующий шаг проверил именно эту более сильную формулировку: main objective
переносится с block-local reconstruction на activation-conditioned
layer-output contributions по реальным captured layer inputs, а exact-program
economics усиливаются отдельными штрафами и explicit decision head для выбора
``atom vs literal''. Это действительно приближает loss к downstream behavior
слоя. Но первая честная реализация показывает, что и этого пока недостаточно.
Exact path остаётся identical by construction, approximate
\texttt{lo\_only} surface не превосходит WAL-DR
(stable \texttt{4}-window gate даёт \texttt{PPL = 2.7463} против
\texttt{2.7325} у \texttt{dr\_only}), а exact program почти полностью
схлопывается в literal fallback: average program length составляет
\texttt{13.0018 / 13.0000} при raw length \texttt{12}, low-level calls лишь
\texttt{0.0140 / 0.0000}. То есть даже objective, привязанный к layer-output
behavior, ещё не создаёт короткую exact ISA, если atom bank не становится
действительно экономически полезным.

\emph{WAL with Learnable High-Expressivity Atoms (WAL-LHA)}.
Следующий шаг оставил WAL-LO objective surface почти без изменений, но сменил
сами low-level atoms: теперь каждый atom задаётся маленьким expressive module,
который производит context-dependent token logits поверх learned base phrase
prior, а training дополнительно включает atom-selection loss и прямой bonus за
использование atoms, когда они улучшают match относительно static bank. Это
даёт первый mixed result в правильную сторону: approximate
\texttt{lha\_only} surface становится лучшей в текущей objective-линии
(stable gate даёт \texttt{PPL = 2.7179}, немного лучше
\texttt{dr\_only = 2.7325} и явно лучше \texttt{lo\_only = 2.7463}), а
sample layer-output relMSE улучшается до примерно \texttt{1.34 / 1.21}. Но
exact economics почти не сдвигаются: average program length остаётся
\texttt{13.0006 / 13.0000} при raw length \texttt{12}, а low-level calls лишь
\texttt{0.0083 / 0.0000}. То есть более выразительные atoms уже помогают как
better approximate language, но всё ещё не становятся economically useful exact
instructions.

\emph{WAL with Strict Budgeted Compatibility (WAL-SBC)}.
Следующий шаг уже не усиливал atoms, а менял сам contract exactness: кроме
старого strict legacy path, теперь существует \texttt{budgeted\_exact} route,
где atom принимается только если одновременно выполняются жёсткие block-level и
activation-conditioned output-level budgets, а fallback идёт сначала в маленький
residual phrase bank и лишь затем в literals. Это важный, но отрицательный
diagnostic result. На первом честном gate уже на \texttt{16} окнах
\texttt{budgeted\_exact} почти не отличается от legacy по quality/runtime
(\texttt{PPL 2.8055 -> 2.8069}, throughput и peak VRAM практически те же), то
есть смена exactness-contract сама по себе не разрушает систему. Но program
economics почти не сдвигаются: average program length остаётся
\texttt{12.9983 / 12.9985} при raw length \texttt{12}, atom calls фактически
нулевые, а residual calls лишь \texttt{0.0017 / 0.0014}. Значит, пересмотреть
definition of exact match действительно было нужно, но первая строгая
budget-геометрия всё ещё слишком консервативна, чтобы materially купить usable
exact-program surface.

Следующий шаг уже был намеренно двухшаговым и percentile-driven: сначала
\texttt{m27\_wal\_sbc\_budget\_profile.py} измеряет реальные CDF для
\texttt{atom\_output\_rel\_mse} и best-residual error на сохранённом
expressive-atom ISA, затем \texttt{m27\_wal\_sbc\_tune\_proto.py} гоняет
только узкий grid с screening на \texttt{4} окнах и честным gate лишь для
лучших \texttt{2} конфигураций на \texttt{16} окнах. Это даёт более чистый
verdict: loose percentile budgets действительно мгновенно оживляют non-literal
path (у лучших finalists \texttt{mean\_nonliteral\_calls \~ 3.85},
\texttt{avg\_program\_length \~ 4.8}), но ломают sacred quality bar
(\texttt{16}-window \texttt{PPL delta vs strict legacy \~ +0.37}). Значит,
этот результат уже оправдывает переход к WAL-CDA: следующий шаг должен менять
сам basis через экономически ограниченный context-aware atom layer, а не делать
ещё один blind widening budgets. При этом сам CDA-probe тоже должен быть жёстко
ограничен: embeddings для дискретных identity-сигналов, только cheap detached
activation summary вместо полного live input, factorized low-rank delta поверх
base prior и soft expected-cost usefulness surrogate против лучшего старого
accepted path на том же budget-contract. Если SBC когда-то и продолжать отдельно, то
только через существенно более tight или shape-aware acceptance geometry.

\item \textbf{M27.} WAL grammar induction: маленький transformer обучается
предсказывать следующий $(s,i)$ токен; embedding-пространство токенов и
per-layer grammaticality становятся новой диагностикой.
\item \textbf{M28.} Component-specific dialects \texttt{WAL-Attn / WAL-MLP /
WAL-Head / WAL-Norm}.
\item \textbf{M29.} WAL $\leftrightarrow$ WAL translation между двумя моделями
одного семейства как proxy weight merging без fine-tuning.
\item \textbf{M30.} Production-ready inference + первая публичная статья и
open release.
\end{itemize}

\section{Что дописывать дальше}

Этот файл дальше нужно дополнять по четырём направлениям:

\begin{itemize}[leftmargin=1.5em]
\item формальная постановка и related work;
\item точное описание encode/decode/codebook pipeline;
\item полная таблица materialized-vs-packed runtime с явным разделением
deployable и research веток;
\item отдельная методологическая вставка о same-encoding compare;
\item расширенный раздел с отрицательными результатами и engineering lessons;
\item единая WAL-таблица для M26 stage-2 probes B0--B3 с честным pass/fail gate
по exactness, hit rate, throughput и peak VRAM.
\end{itemize}